%! $n$th Roots \& Rational Exponents — Unit 6, Lesson 6.1 — Warm-Up (Answer Key)
\documentclass[10pt]{article}
\usepackage{algebra2-article}
\usepackage{algebra2-key}   % pulls in -boxes; do NOT also load -boxes
\newcommand{\TallMath}[1]{$\displaystyle #1\rule[-1.4em]{0pt}{3.2em}$}

\begin{document}

\pageheader{Unit 6, Lesson 6.1}{Warm-Up --- Answer Key}
\namedateperiod

\vspace{0.10in}

\begin{notesbox}{Getting Ready: Skills We'll Use Today}
\small These three skills power today's lesson on $n$th roots and rational exponents. Work quickly.

\vspace{4pt}
\textbf{1. Powers you should recognize on sight.} Evaluate each.
\begin{center}
\renewcommand{\arraystretch}{1.5}
\begin{tabular}{@{}l l l l@{}}
(a) \; $2^5=$ \ans{32} & (b) \; $3^4=$ \ans{81} &
(c) \; $(-3)^3=$ \ans{$-27$} & (d) \; $(-2)^4=$ \ans{16} \\
\end{tabular}
\end{center}
Two of these bases were negative, but only one \emph{answer} came out negative. Which exponents
keep a negative sign, and which destroy it? \ansline{Odd exponents keep the sign (an odd number of
negative factors stays negative); even exponents destroy it (the negatives pair off), so an even
power is never negative.}

\vspace{4pt}
\textbf{2. The exponent properties, with whole-number exponents.} Simplify. Leave your answer in
exponent form.
\begin{center}
\renewcommand{\arraystretch}{1.5}
\begin{tabular}{@{}l l l l@{}}
(a) \; $x^3\cdot x^5=$ \ans{$x^{8}$} & (b) \; $\dfrac{x^7}{x^2}=$ \ans{$x^{5}$} &
(c) \; $\left(x^4\right)^3=$ \ans{$x^{12}$} & (d) \; $x^{-2}=$ \ans{$\dfrac{1}{x^{2}}$} \\
\end{tabular}
\end{center}
In (a) you \ans{add} the exponents; in (c) you \ans{multiply} them.

\vspace{4pt}
\textbf{3. Fraction arithmetic --- you will need it in the exponents themselves.} Simplify.
\begin{center}
\renewcommand{\arraystretch}{1.6}
\begin{tabular}{@{}l l l@{}}
(a) \; $\dfrac{1}{2}+\dfrac{1}{3}=$ \ans{$\dfrac{5}{6}$} &
(b) \; $\dfrac{3}{4}-\dfrac{1}{2}=$ \ans{$\dfrac{1}{4}$} &
(c) \; $\dfrac{2}{3}\cdot 3=$ \ans{$2$} \\
\end{tabular}
\end{center}

\end{notesbox}

\begin{teachernote}
Every item is a piece of today's lesson in disguise --- debrief all three in about four minutes.
\textbf{Item 1} is the perfect-power vocabulary students will read \emph{backwards} all period
($32=2^5$ becomes $\sqrt[5]{32}=2$), and its follow-up question is the algebraic reason Lesson 6.0's
even/odd index rule works: an even power can never be negative, so an even root of a negative
number cannot exist. \textbf{Item 2} is the entire toolkit --- today's claim is that these same four
rules keep working when the exponents become fractions, and item 2(c) in particular is the
derivation of $a^{1/n}$: if $(a^{1/n})^n=a^{n/n}=a$, then $a^{1/n}$ \emph{has} to be the $n$th root.
\textbf{Item 3} is not filler: once exponents are fractions, ``add the exponents'' means adding
fractions, and $\sqrt{x}\cdot\sqrt[3]{x}=x^{1/2+1/3}=x^{5/6}$ is exactly item 3(a). Students who
cannot do item 3 quickly will stall in Tier E, not on the radicals but on the arithmetic --- note
who they are now.
\end{teachernote}

\end{document}
